\documentclass[11pt]{amsart}
\usepackage[T1]{fontenc}
\usepackage{lmodern,amsmath,amssymb,amsthm,microtype}
\usepackage[margin=1.15in]{geometry}
\usepackage[colorlinks=true,linkcolor=blue,citecolor=blue,urlcolor=blue]{hyperref}
% Repository locator for the public (non-anonymous) build of
% non_mf_groups_exist.tex.  Kept in a separate file so that the anonymous
% build has an anonymous source archive and not merely an anonymous PDF:
% exclude this file from a double-blind submission and the manuscript still
% compiles, printing "supplied with this submission as supplementary
% material" in place of the link.
%
% Loaded only when \anonymousfalse is set in the preamble.
\newcommand{\coderepository}{https://github.com/SauersML/group-approximation}
\newcommand{\coderepositoryname}{github.com/SauersML/group-approximation}

\hypersetup{pdftitle={Infinite simple Kazhdan groups that are limits of finite simple groups},pdfauthor={\paperauthorname},pdfkeywords={property (T), sofic groups, hyperlinear groups, LEF groups, simple groups, subshifts, expanders}}
\newtheorem{theorem}{Theorem}
\newtheorem{corollary}[theorem]{Corollary}
\newcommand{\F}{\mathbb F}
\newcommand{\Z}{\mathbb Z}
\newcommand{\LC}{\operatorname{LC}}
\newcommand{\EL}{\operatorname{EL}}
\newcommand{\GL}{\operatorname{GL}}
\newcommand{\SL}{\operatorname{SL}}
\newcommand{\PSL}{\operatorname{PSL}}
\newcommand{\doi}[1]{\href{https://doi.org/#1}{\nolinkurl{doi:#1}}}
\newcommand{\arxiv}[1]{\href{https://arxiv.org/abs/#1}{\nolinkurl{arXiv:#1}}}
\raggedbottom
\emergencystretch=1.5em

\begin{document}
\title{Infinite simple Kazhdan groups that are limits of finite simple groups}
\author{\paperauthorname}
\date{}
\subjclass[2020]{Primary 20E32; Secondary 05C48, 16S35, 20E26, 20F10, 22D55, 37B10, 46L10}
\keywords{Property (T), sofic groups, hyperlinear groups, locally embeddable into finite groups, simple groups, minimal subshifts, expanders}

\begin{abstract}
For every infinite minimal subshift $X$, the group
$G_X=\EL_3(\LC(X,\F_2)\rtimes\Z)$ is an infinite, finitely generated,
simple group with property~\textup{(T)}. It is a limit of finite simple
groups $\SL_{3N}(\F_2)$ in the space of marked groups, and their Cayley
graphs form a family of expanders. So $G_X$ is locally embeddable into
finite groups, sofic and hyperlinear. This answers the question of Brown
and Ozawa whether an infinite simple Kazhdan group can be hyperlinear,
and Pestov's sofic version of it. A finitely generated group is locally
embeddable into finite groups if and only if it is a subgroup of an
infinite finitely generated simple Kazhdan group with this property, and
one such group contains every finitely generated linear group. Every
Turing degree is the word-problem degree of some $G_X$.
\end{abstract}
\maketitle

Can an infinite simple group with property~\textup{(T)} be hyperlinear?
Brown asked this in 2001 as a question about embeddings into unitary
groups of McDuff factors that embed in an ultrapower $\mathcal R^\omega$
of the hyperfinite $\mathrm{II}_1$ factor~\cite[\S11, Question~7]{Brown}.
Ozawa stated the hyperlinear form in 2003~\cite[p.~527]{Ozawa}, and
Pestov's Open question~9.1 adds the sofic form~\cite{Pestov}. Hyperlinear
groups were named by R\u{a}dulescu~\cite{Radulescu}, and sofic groups were
introduced by Gromov~\cite{Gromov} and named by Weiss~\cite{Weiss}. The
groups below answer all three forms positively.

\begin{theorem}\label{thm:main}
Let $X\subseteq A^{\Z}$ be an infinite minimal subshift over a finite
alphabet, with shift $T$, and let $\LC(X,\F_2)$ be the ring of locally
constant functions $X\to\F_2$. Then
\[
  G_X=\EL_3\bigl(\LC(X,\F_2)\rtimes_T\Z\bigr)
\]
is an infinite, finitely generated, simple group with Kazhdan's
property~\textup{(T)}. It is the limit, in the space of marked groups, of
finite simple groups $\SL_{3N}(\F_2)$ whose Cayley graphs, with respect to
the images of a fixed generating set, form a family of expanders. So $G_X$
is locally embeddable into finite groups \textup{(LEF)}, sofic and
hyperlinear. The same holds for $\EL_n\bigl(\LC(X,\F_2)\rtimes_T\Z\bigr)$
for every $n\ge3$, with $\SL_{nN}(\F_2)$ in place of $\SL_{3N}(\F_2)$.
\end{theorem}

Property~\textup{(T)}, introduced by Kazhdan~\cite{Kazhdan}, follows from
the theorem of Ershov and Jaikin-Zapirain~\cite{EJZ}. The finite models
use periodic approximation, as in the proof by Grigorchuk and Medynets
that topological full groups of minimal Cantor systems are
LEF~\cite[Theorem~2.6]{GM}, and the expanders arise as in Kassabov's
construction~\cite{Kassabov}. Simplicity also follows from Stepanov's
theorem on the normal structure of $\GL_n$~\cite[Theorem~4.4]{Stepanov},
since the ring is simple and any two of its elements $p,q$ satisfy
$pr+qs=0$ with $(r,s)\ne0$, by a dimension count on one cylinder. Our
proof is direct: a nontrivial normal subgroup contains a nontrivial
commutator lying in a copy of the finite simple group $\GL_d(\F_2)$ over
a clopen tower, so it contains this group and with it an elementary
matrix. For derived topological full groups, Matui showed in the same way
that a nontrivial normal subgroup meets a simple union of alternating
groups on towers~\cite[Lemma~3.4 and Theorem~4.9]{Matui}. Thom
constructed a finitely generated Kazhdan LEF group that is not residually
finite~\cite{Thom}, but his example is not simple.

The same proof applies to the lamplighter action of a group $\Delta$ on
$\F_2^{\Delta}$. It shows that a finitely generated group is LEF if and
only if it is a subgroup of an infinite finitely generated simple Kazhdan
LEF group (Corollary~\ref{cor:lef}), and one such group contains every
finitely generated linear group.

\section{Proof of Theorem~\ref{thm:main}}

We write the proof for $n=3$. For $n\ge3$ replace $3$ by $n$ throughout,
so that $d=n(2w+1)$ below.

\subsection*{The ring and property (T)}
Write $e_U$ for the indicator of a clopen set $U\subseteq X$. Our
conventions are $(Tx)_t=x_{t+1}$ and
\[
 R=\LC(X,\F_2)\rtimes_T\Z=\bigoplus_{j\in\Z}\LC(X,\F_2)\,u^j,
 \qquad ufu^{-1}=f\circ T^{-1}.
\]
So every element of $R$ is uniquely a finite sum $\sum_jf_ju^j$, the
product is $(fu^i)(f'u^j)=f\,(f'\circ T^{-i})\,u^{i+j}$, and
$ue_Uu^{-1}=e_{TU}$.

The ring $R$ is generated by $u^{\pm1}$ and the letter indicators
$e_a=e_{\{x:x_0=a\}}$, since products of the translates $u^je_au^{-j}$
are the cylinder indicators, which span $\LC(X,\F_2)$. Let
$G=G_X=\EL_3(R)$ be the subgroup of $\GL_3(R)$ generated by the elementary
matrices $e_{ij}(r)=I_3+rE_{ij}$, $i\ne j$, $r\in R$, where $E_{ij}$ is
the usual matrix unit. With $[g,h]=ghg^{-1}h^{-1}$,
\begin{equation}\label{eq:elementary}
 \begin{aligned}
 e_{ij}(r+s)&=e_{ij}(r)e_{ij}(s),\qquad i\ne j,\\
 [e_{il}(r),e_{lj}(s)]&=e_{ij}(rs),\qquad i,j,l\text{ distinct}.
 \end{aligned}
\end{equation}
So the matrices $e_{ij}(s)$ with $s\in\{u,u^{-1}\}\cup\{e_a:a\in A\}$
generate $G$, as $1=\sum_ae_a$. By Ershov and
Jaikin-Zapirain~\cite[Theorem~1.1]{EJZ}, $\EL_n$ of every finitely
generated associative unital ring has property~\textup{(T)} for
$n\ge3$. So $G$ has property~\textup{(T)}. It is infinite because
$e_{12}(\LC(X,\F_2))$ is infinite.

\subsection*{Finite models}
A group is LEF~\cite{VershikGordon} if every finite subset embeds
injectively into a finite group, preserving all products that stay in
that subset.

Fix $x\in X$ and $\ell\ge0$. By minimality the forward orbit of $x$ is
dense, so every word of $X$ occurs in $x_{[0,\infty)}$, and
$x_{[0,2\ell)}$ recurs at arbitrarily large positions. Choose such a
position $m$ with all words of length $2\ell+1$ of $X$ occurring in
$x_{[0,m+2\ell)}$. Since $x_{[m,m+2\ell)}=x_{[0,2\ell)}$, the $m$-periodic
sequence $y_\ell$ that agrees with $x$ on $[0,m)$ agrees with $x$ on
$[0,m+2\ell)$. Every window of length $2\ell+1$ of $y_\ell$ is a
translate, by a multiple of $m$, of a window starting in $[0,m)$, and
every word of $X$ of this length starts at some position in $[0,m)$. So
$y_\ell$ has the same words of length $2\ell+1$ as $X$. Let $N_\ell$ be
the least period of $y_\ell$. Then $N_\ell$ is at least the number of
these words, which tends to infinity with $\ell$, since a subshift with
boundedly many words of each length is finite~\cite{MorseHedlund38}.

On $\F_2^{\Z/N_\ell\Z}$ let $P\delta_t=\delta_{t+1}$. For $f$ depending
only on coordinates in $[-\ell,\ell]$ let
$D_\ell(f)\delta_t=f(T^ty_\ell)\delta_t$, where $f(T^ty_\ell)$ is the value
of $f$ at any point of $X$ that agrees with $T^ty_\ell$ on $[-\ell,\ell]$.
When every $f_j$ depends only on coordinates in $[-\ell,\ell]$, put
$\varphi_\ell(\sum_jf_ju^j)=\sum_jD_\ell(f_j)P^j$. Since
$PD_\ell(f)P^{-1}=D_\ell(f\circ T^{-1})$, for fixed $r,s\in R$ the
identities $\varphi_\ell(r+s)=\varphi_\ell(r)+\varphi_\ell(s)$ and
$\varphi_\ell(rs)=\varphi_\ell(r)\varphi_\ell(s)$ hold for large $\ell$.
If $r=\sum_jf_ju^j\ne0$, then for large $\ell$ some $D_\ell(f_j)\ne0$, and
the terms $D_\ell(f_j)P^j$ have disjoint supports once $N_\ell>2|j|$ for
all exponents $j$ of $r$, so $\varphi_\ell(r)\ne0$.

Let $F=\F_2\langle\tau_+,\tau_-,\tau_a:a\in A\rangle$ be the free algebra,
and let $\pi\colon F\to R$ and $\rho_\ell\colon F\to M_{N_\ell}(\F_2)$ be
the ring homomorphisms with $\pi(\tau_\pm)=u^{\pm1}$, $\pi(\tau_a)=e_a$,
$\rho_\ell(\tau_\pm)=P^{\pm1}$ and $\rho_\ell(\tau_a)=D_\ell(e_a)$. Both
are onto. For $\rho_\ell$, the shifts $T^ty_\ell$, $0\le t<N_\ell$, are
distinct, so products of the matrices $P^iD_\ell(e_a)P^{-i}$ give every
diagonal matrix unit, and with $P$ every matrix unit. So $G$ and
$\EL_3(M_{N_\ell}(\F_2))=\SL_{3N_\ell}(\F_2)$ are quotients of $\EL_3(F)$,
which is generated by the matrices $e_{ij}(1)$ and $e_{ij}(\tau)$ with
$\tau\in\{\tau_\pm,\tau_a\}$. The equality holds since transvections
generate $\SL_{3N_\ell}(\F_2)$, and those inside one block are commutators
of those between blocks. A word of length $\lambda$ in the $e_{ij}(\tau)$
has entries of degree at most $\lambda$ in $F$, and
$\varphi_\ell\circ\pi=\rho_\ell$ on these entries for large $\ell$.
Applied to the entries of the word minus $I_3$, the previous paragraph
shows that for large $\ell$ the word is trivial in $G$ if and only if it
is trivial in $\SL_{3N_\ell}(\F_2)$. So these finite simple groups
converge to $G$ in the space of marked
groups~\cite{Grigorchuk,Champetier}, and $G$ is LEF. Since $\EL_3(F)$ has
property~\textup{(T)}~\cite[Theorem~1.1]{EJZ}, its finite quotients form
a family of expanders~\cite{Margulis,Kassabov}. In both quotients
$e_{ij}(1)$ is the image of $\prod_ae_{ij}(\tau_a)$, as $\sum_ae_a=1$ and
$\sum_aD_\ell(e_a)=I$, so this holds for the images of the
$e_{ij}(\tau)$. LEF groups are sofic~\cite[Example~4.5]{Pestov}, and
sofic groups are hyperlinear~\cite[Theorem~2]{ElekSzabo}.

\subsection*{Simplicity}
Let $1\ne K\trianglelefteq G$ and $1\ne g\in K$, and let $w\ge0$ bound
the absolute values of the exponents in the entries of $g$ and $g^{-1}$.
Call a clopen set $V$ small if $V\cap T^tV=\varnothing$ for
$0<|t|\le2w$ and every $f\circ T^t$, with $|t|\le w$ and $f$ a
coefficient of an entry of $g$ or $g^{-1}$, is constant on $V$.
Minimality and infiniteness imply that $T$ has no periodic points, so
every nonempty clopen set contains a small one.

Some $h=e_{ij}(e_V)$ with $V$ small does not commute with $g$.
Indeed, $g$ commutes with $e_{ij}(r)$ exactly when $g_{pi}r=0$ for
$p\ne i$, $rg_{jq}=0$ for $q\ne j$, and $g_{ii}r=rg_{jj}$. For
$c=\sum_tc_tu^t$ and small $V$ we have $ce_V=\sum_tc_te_{T^tV}u^t$ and
$e_Vc=\sum_te_Vc_tu^t$. If $c_t\ne0$, it equals $1$ on a nonempty clopen
set $U$, and a small set $V\subseteq T^{-t}U$ gives $ce_V\ne0$, while a
small set $V\subseteq U$ gives $e_Vc\ne0$. If $g$ commuted with every
such $h$, then $g_{pi}e_V=0$ for all small $V$ and all $p\ne i$, so $g$
would be diagonal. In $g_{ii}e_V=e_Vg_{jj}$ the coefficients at $u^t$
with $t\ne0$ are supported in the disjoint sets $T^tV$ and $V$, so they
vanish, and as above the coefficients of $g_{ii}$ and $g_{jj}$ at these
$u^t$ vanish. The constant coefficients agree on every small $V$, hence
everywhere, since every nonempty clopen set contains a small one. So
$g=cI_3$ with $c\in\LC(X,\F_2)$, and comparing constant coefficients in
$cc^{-1}=1$ gives $c=1$ and $g=1$, a contradiction.

Fix such $h$ and $V$, put $d=3(2w+1)$, and for $|a|,|b|\le w$ put
$\epsilon_{ab}=e_{T^aV}u^{a-b}$. They are nonzero, as $h\ne I_3$ gives
$V\ne\varnothing$. Since
$\epsilon_{ab}\epsilon_{a'b'}=e_{T^aV\cap T^{a-b+a'}V}\,u^{a-b+a'-b'}$ and
$V\cap T^tV=\varnothing$ for $0<|t|\le2w$, we get
$\epsilon_{ab}\epsilon_{a'b'}=\delta_{ba'}\epsilon_{ab'}$. Index the
coordinates of $\F_2^d$ by pairs $(p,a)$ with $1\le p\le3$ and
$|a|\le w$. Then $\psi(E_{(p,a),(q,b)})=\epsilon_{ab}E_{pq}$ defines an
injective multiplicative linear map $\psi\colon M_d(\F_2)\to M_3(R)$, and
$M\mapsto I_3-\psi(I_d)+\psi(M)$ embeds $\GL_d(\F_2)$ in $\GL_3(R)$. Its
image $H$ lies in $G$. Indeed, transvections generate $\GL_d(\F_2)$. The
transvection between $(p,a)$ and $(q,b)$ maps to $e_{pq}(\epsilon_{ab})$
if $p\ne q$. If $p=q$, it is the commutator of the transvections between
$(p,a)$ and $(p',a)$ and between $(p',a)$ and $(p,b)$, where $p'\ne p$.

Put $k=[g,h]\in K\setminus\{1\}$. If $|a|,|b|\le w$ and $f,f'$ are
coefficients as above, then
\[
  fu^a\,e_V\,f'u^b=f\,e_{T^aV}\,(f'\circ T^{-a})\,u^{a+b}
  \in\{0,\epsilon_{a,-b}\},
\]
since $f\circ T^a$ and $f'$ are constant on $V$. The entries of
$ghg^{-1}-I_3=g\,e_VE_{ij}\,g^{-1}$ are sums of such products, and
$h^{-1}=h=I_3+\epsilon_{00}E_{ij}$. The image $\psi(M_d(\F_2))$ is the set
of matrices with entries in the span of the $\epsilon_{ab}$, and it is
closed under products. So $k-I_3=(ghg^{-1}-h)h$ and
$k^{-1}-I_3=h(ghg^{-1}-h)$ lie in it, say $k=I_3+\psi(M)$ and
$k^{-1}=I_3+\psi(M')$. As $\psi$ is injective and multiplicative,
$(I_d+M)(I_d+M')=I_d$, so $k$ is the image of $I_d+M$ and $k\in H$. Then
$K\cap H$ is a nontrivial normal subgroup of
$H\cong\GL_d(\F_2)=\PSL_d(\F_2)$, which is simple as $d\ge3$. So
$H\subseteq K$, and $e_{pq}(e_V)\in K$ for all $p\ne q$.

Finally, $J=\{r\in R:e_{pq}(r)\in K\text{ for all }p\ne q\}$ is a
two-sided ideal. It is additive by~\eqref{eq:elementary}, and for
$l\notin\{p,q\}$ we have $e_{pq}(sr)=[e_{pl}(s),e_{lq}(r)]$ and
$e_{pq}(rs)=[e_{pl}(r),e_{lq}(s)]$. It contains $e_V$, so it contains
every $e_{T^aV}=u^ae_Vu^{-a}$. By minimality and compactness finitely
many of these sets cover $X$, so
$1=1-\prod_i(1-e_{T^{a_i}V})\in J$. So $J=R$ and $K=G$.\qed

\subsection*{Brown's formulation}
As $G$ is hyperlinear, $L(G)$ embeds in
$\mathcal R^\omega$~\cite[Proposition~7.1]{Ozawa}. Since $G$ is infinite
and simple, its nontrivial conjugacy classes are infinite and $G$ is not
residually finite. So $L(G)$ is a $\mathrm{II}_1$ factor, and
$L(G)\mathbin{\bar\otimes}\mathcal R$ is a McDuff factor~\cite{McDuff}
that embeds in $\mathcal R^\omega$ and whose unitary group contains $G$.
This is Brown's formulation. Kirchberg proved that a Kazhdan group with
the factorization property is residually
finite~\cite[Theorem~1.1]{Kirchberg}. So $G$ does not have the
factorization property, and $C^*(G)$ does not have the local lifting
property~\cite[p.~527]{Ozawa}.

\section{LEF groups}

Every countable group embeds in a finitely generated simple
group~\cite{Gorjuskin,Schupp}, and Kionke and Schesler proved that every
finitely generated residually finite group embeds in a finitely generated
simple LEF group~\cite[Theorem~1.2]{KionkeSchesler}. The next corollary
adds property~\textup{(T)} and applies to every LEF group. Since LEF
passes to subgroups, it characterizes LEF groups.

\begin{corollary}\label{cor:lef}
A finitely generated group is LEF if and only if it is a subgroup of an
infinite finitely generated simple Kazhdan LEF group.
\end{corollary}

\begin{proof}
Subgroups of LEF groups are LEF. Conversely, let $\Gamma$ be LEF and
generated by a finite set $S$, with finite models
$\varphi_n\colon B_n\to Q_n$ on word balls. Then
$\gamma\mapsto(\varphi_n(\gamma))_n$ embeds $\Gamma$ in the algebraic
ultraproduct $\prod_\omega Q_n$, and the doubled regular actions embed
this group in $P=\prod_\omega\operatorname{Sym}(Q_n\times\{1,2\})$ by even
permutations. Every even permutation of a finite set is a commutator of
two permutations~\cite[Theorem~1]{Ore}, so each $s\in S$ is a commutator
of two elements of $P$. Let $\Delta_0\le P$ be generated by $S$ and these
elements, and put $\Delta=\Delta_0\times\Z$. A countable subgroup of an
algebraic ultraproduct of finite groups is LEF, so $\Delta$ is an
infinite finitely generated LEF group with $\Gamma\le[\Delta,\Delta]$.

Let $\Delta$ act on $Z=\F_2^{\Delta}$ by $(\delta x)(h)=x(\delta^{-1}h)$,
and let $L$ be the group of maps $x\mapsto\delta x+f$ with
$\delta\in\Delta$ and $f$ finitely supported. It is generated by $\Delta$
and the map adding $1$ at the coordinate $e$. Every orbit contains
$x+\bigoplus_\Delta\F_2$, so the action is minimal. It is topologically
free: for $f\ne0$ the map $x\mapsto x+f$ has no fixed point, and if
$\delta\ne e$, every cylinder over a finite set $W$ contains points with
$x(h)\ne x(\delta^{-1}h)$ for some
$h\notin W\cup\delta W\cup\operatorname{supp}f$, and $x\mapsto\delta x+f$
moves them. Put $R=\LC(Z,\F_2)\rtimes L$ and $G=\EL_3(R)$. The ring $R$
is generated by the units $u_s^{\pm1}$ for the generators $s$ of $L$ and
the indicator of $\{x:x(e)=1\}$, so $G$ is finitely generated and has
property~\textup{(T)}~\cite[Theorem~1.1]{EJZ}, and it is infinite. Points
with trivial stabilizer in a given ball are dense, so every nonempty
clopen set contains a small one, in the sense of Section~1 with $u_\ell$
for $u^t$, $\ell V$ for $T^tV$, products in $L$ for sums of exponents,
word length for $|t|$, and the ball $B_w$ of $L$ for $\{a:|a|\le w\}$, so
that $d=3|B_w|$. With these replacements the proof of
simplicity in Section~1 applies word for word, so $G$ is simple.

For finite models $\psi_n\colon B_n\to Q'_n$ of $\Delta$ put
$E_n=\F_2^{Q'_n}\times Q'_n$, and attach to $(c,q)\in E_n$ the
configuration $h\mapsto c(q\psi_n(h))$ on $B_n$. Let $x\mapsto\delta x+f$
act on $E_n$ by
$(c,q)\mapsto\bigl(c+\sum_{h\in\operatorname{supp}f}1_{q'\psi_n(h)},q'\bigr)$
with $q'=q\psi_n(\delta^{-1})$. For fixed elements of $L$ and large $n$,
these permutations multiply as in $L$ and move the attached
configurations near $e$ as $L$ moves points of $Z$. Every finite pattern
occurs in $Z$, so a function depending on coordinates in a ball defines a
diagonal matrix $D_n(f)$ by its values at the attached configurations,
and $\sum_\ell f_\ell u_\ell\mapsto\sum_\ell D_n(f_\ell)P_\ell$, with
$P_\ell$ the permutation matrices, is a unital ring homomorphism from $R$
to the algebraic ultraproduct $\prod_\omega M_{|E_n|}(\F_2)$. It maps $G$
to the ultraproduct of the groups $\GL_{3|E_n|}(\F_2)$. The kernel is a
normal subgroup of the simple group $G$ that does not contain
$e_{12}(1)$, so $G$ embeds there and is LEF.

Finally, $\ell\mapsto u_\ell$ embeds $L$ in $\GL_1(R)$. Over $\F_2$, for
units $a,b,c$ of $R$,
\[
 \begin{aligned}
 e_{12}(c)e_{21}(c^{-1})e_{12}(c)\,e_{12}(1)e_{21}(1)e_{12}(1)
 &=\operatorname{diag}(c,c^{-1},1),\\
 \operatorname{diag}(a,a^{-1},1)\operatorname{diag}(b,b^{-1},1)
 \operatorname{diag}\bigl((ba)^{-1},ba,1\bigr)
 &=\operatorname{diag}(aba^{-1}b^{-1},1,1).
 \end{aligned}
\]
So $\ell\mapsto\operatorname{diag}(u_\ell,1,1)$ is an injective
homomorphism $L\to\GL_3(R)$ that maps $[L,L]$ into $G$, and
$\Gamma\le[\Delta,\Delta]\le[L,L]$.
\end{proof}

\begin{corollary}
One infinite finitely generated simple Kazhdan LEF group contains every
finitely presented residually finite group and every finitely generated
linear group.
\end{corollary}

\begin{proof}
Finitely generated linear groups are residually finite by Malcev's
theorem, and up to isomorphism there are countably many groups of either
kind. So their direct sum is a countable residually finite group. Wilson
proved that every countable residually finite group embeds in a
$2$-generator residually finite group~\cite[Theorem~A]{Wilson}. Such a
group is LEF, and Corollary~\ref{cor:lef} applies.
\end{proof}

\section{Word problems}

\begin{corollary}
The word problem of $G_X$ has the Turing degree of the language $L(X)$
of $X$. In particular $G_X$ has solvable word problem if and only if
$L(X)$ is recursive, as for the Fibonacci subshift. Every Turing degree
occurs, so there are continuum many pairwise nonisomorphic groups $G_X$.
\end{corollary}

\begin{proof}
Multiplying out a word in the generators in $\LC(A^{\Z},\F_2)\rtimes\Z$,
which maps onto $R$, gives a matrix with entries $\sum_jf_ju^j$, each
$f_j$ given by a table on the words of some length, using
$uf=(f\circ T^{-1})u$. The word is trivial in $G_X$ if and only if the
tables of its difference from $I_3$ vanish on $L(X)$, so $L(X)$ computes
the word problem. Conversely, by~\eqref{eq:elementary} a word for
$e_{12}(\prod_{t<m}u^{-t}e_{v_t}u^t)$ is computable from
$v=v_0\cdots v_{m-1}$, and it is trivial if and only if $v\notin L(X)$.

For irrational $\alpha\in(0,1)$ let $X_\alpha$ be the infinite minimal
Sturmian subshift, the closure of the set of codings $c(\theta)$,
$\theta\in[0,1)$, with $c(\theta)_t=1$ if and only if
$\theta+t\alpha\bmod1\ge1-\alpha$~\cite{MorseHedlund,Hedlund44},
\cite[Chapter~2]{Lothaire}. So $c(\theta)_t=1$ if and only if $\theta$
lies in the arc $[-(t+1)\alpha,-t\alpha)$ modulo $1$. Passing to the
closure adds no words, so the words of length $m$ of $X_\alpha$ are the
constant values of $c(\theta)_{[0,m)}$ on the arcs between the points
$-j\alpha\bmod1$, $0\le j\le m$, which are distinct as $\alpha$ is
irrational. So $\alpha$ computes $L(X_\alpha)$. Since
$c(\theta)_t=\lfloor\theta+(t+1)\alpha\rfloor-\lfloor\theta+t\alpha\rfloor$,
the word $c(\theta)_{[0,m)}$ has $\lfloor\theta+m\alpha\rfloor$ ones,
within $1$ of $m\alpha$, so $L(X_\alpha)$ computes $\alpha$. Every set
$S\subseteq\mathbb N$ has the degree of the irrational number
$[0;1+\chi_S(0),1+\chi_S(1),\dots]$, since each digit is decided by one
strict comparison of this number with a rational computed from the
earlier digits. There are continuum many degrees, and the degree of the
word problem is an isomorphism invariant.
\end{proof}

For derived topological full groups, Grigorchuk and Medynets proved that
the word problem is decidable if and only if $L(X)$ is
recursive~\cite[Theorem~1.1(3)]{GMpres}.

\section{Questions}

A finitely presented LEF group is residually finite~\cite{VershikGordon},
and the infinite simple group $G_X$ is not, so $G_X$ is not finitely
presented. Finitely presented infinite simple groups with
property~\textup{(T)} exist~\cite{CapraceRemy}. Is there one that is
sofic, or at least hyperlinear? A positive answer would also answer Open
problem~6.1 of Alekseev and Thom~\cite{AlekseevThom}, which asks for
finitely presented sofic groups with property~\textup{(T)} that are not
residually finite.

If $(X,T)$ is topologically conjugate to $(Y,S)$ or to $(Y,S^{-1})$, then
$G_X\cong G_Y$. For topological full groups the converse
holds~\cite{GPS99,BezuglyiMedynets}. Does $G_X\cong G_Y$ imply that $X$ and
$Y$ are flip conjugate, or at least strongly orbit equivalent~\cite{GPS}?

\section*{Origin and authorship}
Claude (Anthropic) found the construction and original proofs under
the author's direction on September~12, 2026, and wrote the manuscript.
The author is responsible for the final manuscript. Proof records are
in the \href{https://github.com/SauersML/group-approximation}{project repository}.

\begin{thebibliography}{99}


\bibitem{AlekseevThom}
V.~Alekseev and A.~Thom,
\emph{Centralizers of sofic approximations of Kazhdan groups},
\arxiv{2608.05362} (2026).

\bibitem{BezuglyiMedynets}
S.~Bezuglyi and K.~Medynets,
\emph{Full groups, flip conjugacy, and orbit equivalence of Cantor minimal
systems}, Colloq. Math. \textbf{110} (2008), 409--429.
\doi{10.4064/cm110-2-6}.

\bibitem{Brown}
N.~P. Brown, \emph{Tracial invariants, classification and
$\mathrm{II}_1$ factor representations of Popa algebras},
\arxiv{math/0111286} (2001).

\bibitem{CapraceRemy}
P.-E. Caprace and B.~R\'emy, \emph{Simplicity and superrigidity of twin
building lattices}, Invent. Math. \textbf{176} (2009), 169--221.
\doi{10.1007/s00222-008-0162-6}.

\bibitem{Champetier}
C.~Champetier, \emph{L'espace des groupes de type fini}, Topology
\textbf{39} (2000), 657--680. \doi{10.1016/S0040-9383(98)00063-9}.

\bibitem{ElekSzabo}
G.~Elek and E.~Szab\'o, \emph{Hyperlinearity, essentially free actions and
$L^2$-invariants. The sofic property}, Math. Ann. \textbf{332} (2005),
421--441. \doi{10.1007/s00208-005-0640-8}.

\bibitem{EJZ}
M.~Ershov and A.~Jaikin-Zapirain,
\emph{Property~\textup{(T)} for noncommutative universal lattices},
Invent. Math. \textbf{179} (2010), 303--347.
\doi{10.1007/s00222-009-0218-2}.

\bibitem{GPS}
T.~Giordano, I.~F. Putnam, and C.~F. Skau,
\emph{Topological orbit equivalence and $C^*$-crossed products},
J. Reine Angew. Math. \textbf{469} (1995), 51--112.
\doi{10.1515/crll.1995.469.51}.

\bibitem{GPS99}
T.~Giordano, I.~F. Putnam, and C.~F. Skau,
\emph{Full groups of Cantor minimal systems},
Israel J. Math. \textbf{111} (1999), 285--320.
\doi{10.1007/BF02810689}.

\bibitem{Gorjuskin}
A.~P. Gorju\v{s}kin, \emph{Imbedding of countable groups in
$2$-generated simple groups}, Math. Notes \textbf{16} (1974), 725--727.
\doi{10.1007/BF01105577}.

\bibitem{Grigorchuk}
R.~I. Grigorchuk, \emph{Degrees of growth of finitely generated groups,
and the theory of invariant means}, Math. USSR-Izv. \textbf{25} (1985),
259--300. \doi{10.1070/IM1985v025n02ABEH001281}.

\bibitem{GM}
R.~Grigorchuk and K.~Medynets,
\emph{On algebraic properties of topological full groups},
Sb. Math. \textbf{205} (2014), 843--861.
\doi{10.1070/SM2014v205n06ABEH004400}.

\bibitem{GMpres}
R.~Grigorchuk and K.~Medynets,
\emph{Presentations of topological full groups by generators and
relations}, J. Algebra \textbf{500} (2018), 46--68.
\doi{10.1016/j.jalgebra.2016.10.027}.

\bibitem{Gromov}
M.~Gromov, \emph{Endomorphisms of symbolic algebraic varieties},
J. Eur. Math. Soc. \textbf{1} (1999), 109--197.
\doi{10.1007/PL00011162}.

\bibitem{Hedlund44}
G.~A. Hedlund, \emph{Sturmian minimal sets},
Amer. J. Math. \textbf{66} (1944), 605--620.
\doi{10.2307/2371769}.

\bibitem{Kassabov}
M.~Kassabov, \emph{Universal lattices and unbounded rank expanders},
Invent. Math. \textbf{170} (2007), 297--326.
\doi{10.1007/s00222-007-0064-z}.

\bibitem{Kazhdan}
D.~A. Kazhdan, \emph{Connection of the dual space of a group with the
structure of its closed subgroups}, Funct. Anal. Appl. \textbf{1} (1967),
63--65. \doi{10.1007/BF01075866}.

\bibitem{KionkeSchesler}
S.~Kionke and E.~Schesler,
\emph{From telescopes to frames and simple groups},
J. Comb. Algebra (2024). \doi{10.4171/JCA/103}.
\arxiv{2304.09307}.

\bibitem{Kirchberg}
E.~Kirchberg, \emph{Discrete groups with Kazhdan's property~$T$ and
factorization property are residually finite},
Math. Ann. \textbf{299} (1994), 551--563.
\doi{10.1007/BF01459798}.

\bibitem{Lothaire}
M.~Lothaire, \emph{Algebraic combinatorics on words}, Encyclopedia Math.
Appl. \textbf{90}, Cambridge Univ. Press, Cambridge, 2002.
\doi{10.1017/CBO9781107326019}.

\bibitem{Margulis}
G.~A. Margulis, \emph{Explicit construction of a concentrator}
(Russian), Problemy Peredachi Informatsii \textbf{9} (1973), no.~4,
71--80.

\bibitem{Matui}
H.~Matui, \emph{Some remarks on topological full groups of Cantor
minimal systems}, Internat. J. Math. \textbf{17} (2006), 231--251.
\doi{10.1142/S0129167X06003448}.

\bibitem{McDuff}
D.~McDuff, \emph{Central sequences and the hyperfinite factor},
Proc. London Math. Soc. (3) \textbf{21} (1970), 443--461.
\doi{10.1112/plms/s3-21.3.443}.

\bibitem{MorseHedlund38}
M.~Morse and G.~A. Hedlund, \emph{Symbolic dynamics},
Amer. J. Math. \textbf{60} (1938), 815--866.
\doi{10.2307/2371264}.

\bibitem{MorseHedlund}
M.~Morse and G.~A. Hedlund,
\emph{Symbolic dynamics II. Sturmian trajectories},
Amer. J. Math. \textbf{62} (1940), 1--42.
\doi{10.2307/2371431}.

\bibitem{Ore}
O.~Ore, \emph{Some remarks on commutators},
Proc. Amer. Math. Soc. \textbf{2} (1951), 307--314.
\doi{10.1090/S0002-9939-1951-0040298-4}.

\bibitem{Ozawa}
N.~Ozawa, \emph{About the QWEP conjecture},
Internat. J. Math. \textbf{15} (2004), 501--530; preprint \arxiv{math/0306067} (2003).
\doi{10.1142/S0129167X04002417}.

\bibitem{Pestov}
V.~G. Pestov, \emph{Hyperlinear and sofic groups: a brief guide},
Bull. Symbolic Logic \textbf{14} (2008), 449--480.
\doi{10.2178/bsl/1231081461}.

\bibitem{Radulescu}
F.~R\u{a}dulescu, \emph{The von Neumann algebra of the non-residually
finite Baumslag group $\langle a,b\mid ab^3a^{-1}=b^2\rangle$ embeds into
$R^\omega$}, \arxiv{math/0004172}
(2000).

\bibitem{Schupp}
P.~E. Schupp, \emph{Embeddings into simple groups},
J. London Math. Soc. (2) \textbf{13} (1976), 90--94.
\doi{10.1112/jlms/s2-13.1.90}.

\bibitem{Stepanov}
A.~V. Stepanov, \emph{On the normal structure of the general linear
group over a ring}, Zap. Nauchn. Sem. S.-Peterburg. Otdel. Mat. Inst. Steklov. (POMI) \textbf{236} (1997), 166--182;
English transl., J. Math. Sci. \textbf{95} (1999), 2146--2155.
\doi{10.1007/BF02169976}.

\bibitem{Thom}
A.~Thom, \emph{Examples of hyperlinear groups without factorization
property}, Groups Geom. Dyn. \textbf{4} (2010), 195--208.
\doi{10.4171/GGD/80}.

\bibitem{VershikGordon}
A.~M. Vershik and E.~I. Gordon,
\emph{Groups that are locally embeddable in the class of finite groups},
Algebra i Analiz \textbf{9} (1997), no.~1, 71--97; English transl.,
St. Petersburg Math. J. \textbf{9} (1998), no.~1, 49--67.

\bibitem{Weiss}
B.~Weiss, \emph{Sofic groups and dynamical systems}, Sankhy\={a} Ser. A
\textbf{62} (2000), 350--359.

\bibitem{Wilson}
J.~S. Wilson, \emph{Embedding theorems for residually finite groups},
Math. Z. \textbf{174} (1980), 149--157.
\doi{10.1007/BF01293535}.


\end{thebibliography}
\end{document}
